% Candidate rules for payment.  Written by hand; see docs/one-order-coverage.md.

\section{Payment}

\begin{definition}[Payment time limit (C-11)]
  \label{def:c11-payment-time-limit}
  \uses{def:c08-delivery-status, def:instant}
  \notready
  A confirmed service waits for payment until the payment time limit the
  airline gave the seller.

  Open point: the source does not say what happens to the service when the
  limit passes without payment.

  Source: \href{https://standards.atlassian.net/wiki/spaces/EASD/pages/574849062/Order}{EASD guide, Order}.
\end{definition}

\begin{definition}[Committing to pay nothing (C-12)]
  \label{def:c12-zero-commitment}
  \uses{def:c08-delivery-status, def:c01-structure}
  \notready
  An order item with nothing to pay goes ahead when the seller commits to pay
  zero.

  Open point: the source does not say which delivery status the item's
  services move to.

  Source: \href{https://standards.atlassian.net/wiki/spaces/EASD/pages/2593488900/Commitment+to+pay+when+no+amount+due}{EASD guide, Commitment to pay when no amount due}.
\end{definition}

\begin{definition}[Paying with several methods (C-13)]
  \label{def:c13-several-methods}
  \uses{def:payment}
  Several payment methods can pay for one order item in a single payment.  The
  payment succeeds only if every method succeeds; if any fails, the whole
  payment fails.

  Source: \href{https://standards.atlassian.net/wiki/spaces/EASD/pages/2592899087/Use+of+multiple+payment+methods+to+pay+for+a+single+item}{EASD guide, Use of multiple payment methods}.
\end{definition}

\begin{theorem}[Amount charged across methods (C-14)]
  \label{thm:c14-amount-covers-price}
  \uses{def:c13-several-methods, thm:c04-item-unsplit}
  \notready
  The amounts charged across the methods cover the item's price plus any
  surcharges.

  Open point: the source's example adds a card surcharge on top of the price
  and states no rule, so it is unclear whether an item paid in full means the
  price was paid or the price plus surcharges.

  Source: \href{https://standards.atlassian.net/wiki/spaces/EASD/pages/2592899087/Use+of+multiple+payment+methods+to+pay+for+a+single+item}{EASD guide, Use of multiple payment methods}.
\end{theorem}

\begin{definition}[Refunds across payment methods (C-17)]
  \label{def:c17-refund-methods}
  \uses{def:c13-several-methods, thm:c14-amount-covers-price, def:accept-refund, thm:refunds-le-payments}
  \notready
  A refund goes back to the payment methods that paid.

  Open point: the source does not say how a refund is split across several
  methods, or whether surcharges are refunded.  If surcharges are refunded, a
  refund can exceed the item's price, which bears on the open question about
  refunds and payments.

  Sources: \href{https://standards.atlassian.net/wiki/spaces/EASD/pages/2592899087/Use+of+multiple+payment+methods+to+pay+for+a+single+item}{EASD guide, Use of multiple payment methods};
  \href{https://standards.atlassian.net/wiki/spaces/EASD/pages/574783569/Cancellation+of+Order+Items}{EASD guide, Cancellation of Order Items}.
\end{definition}

% Second survey round: candidates written from the survey notes in docs/one-order-survey.md.

\begin{definition}[Paid order item (C-126)]
  \label{def:payment-paid-order-item}
  \uses{def:payment, def:c01-structure, def:c08-delivery-status, thm:c14-amount-covers-price, def:c18-change}
  \notready
  An order item is paid once the successful payments allocated to it,
  together with any value carried over from the item it replaced, cover its
  amount due. A \texttt{PENDING} payment does not count. After a change, the
  original payment stays on record and its allocation moves to the replacing
  item.

  Case to check: In the change worked example the new item costs 910.00 and
  the only payment made after the change is 60.00. The item is paid only if
  the 1000.00 payment made for the old item counts toward it.

  Open point: The source speaks of unpaid, partly paid and paid items but
  never states the test. The business must decide whether surcharges count
  toward the amount due and how value carried over from a replaced item is
  counted.

  Sources: \href{https://standards.atlassian.net/wiki/spaces/EASD/pages/574849062/Order}{EASD guide, Order};
  \href{https://retailing.iata.org/armi/docs/PAYORD/}{IATA capability PAYORD};
  \href{https://standards.atlassian.net/wiki/spaces/EASD/pages/842301447/Customer+requests+a+change+of+a+paid+order+prior+to+travel+resulting+in+additional+collection+and+a+refund}{EASD guide, Customer requests a change of a paid order prior to travel resulting in additional collection and a refund};
  \href{https://standards.atlassian.net/wiki/spaces/ARMI/pages/584318980/General+Guidance}{ARMI, General Guidance}.
\end{definition}

\begin{definition}[Payment status and its moves (C-127)]
  \label{def:payment-payment-status}
  \uses{def:payment, def:instant}
  \notready
  A payment transaction is \texttt{PENDING}, \texttt{SUCCESSFUL} or
  \texttt{FAILED}. \texttt{SUCCESSFUL} means the airline is certain to
  collect the funds, \texttt{FAILED} means it knows it will not, and
  \texttt{PENDING} means it cannot yet tell; a redirected payment starts
  \texttt{PENDING} and later moves to \texttt{SUCCESSFUL} or
  \texttt{FAILED}. The moment a payment becomes \texttt{SUCCESSFUL} is its
  payment commitment time.

  Case to check: A settlement plan payment is \texttt{SUCCESSFUL} at order
  time, then the seller remits only 75 percent and the platform marks the
  clearances Not Settled or claws them back. The order still shows the
  payment as \texttt{SUCCESSFUL}.

  Open point: The source does not say whether \texttt{SUCCESSFUL} and
  \texttt{FAILED} are final. The business must decide whether a
  \texttt{SUCCESSFUL} payment can later fail, for example after a card
  chargeback or a seller's default on a settlement plan, and what then
  happens to the items it paid.

  Sources: \href{https://standards.atlassian.net/wiki/spaces/EASD/pages/574849161/Payment+Codes}{EASD guide, Payment Codes};
  \href{https://standards.atlassian.net/wiki/spaces/EASD/pages/574423154/Payment+Status}{EASD guide, Payment Status};
  \href{https://standards.atlassian.net/wiki/spaces/EASD/pages/574455927/Pay+Offline+Using+Redirection}{EASD guide, Pay Offline Using Redirection};
  \href{https://standards.atlassian.net/wiki/spaces/SWO/pages/574653091/Enhanced+and+Simplified+Distribution}{Settlement with Orders, Enhanced and Simplified Distribution};
  \href{https://standards.atlassian.net/wiki/spaces/SWO/pages/574718145/Remittance+of+Funds}{Settlement with Orders, Remittance of Funds}.
\end{definition}

\begin{theorem}[Effects of a payment outcome (C-128)]
  \label{thm:payment-payment-outcome-effects}
  \uses{def:payment-payment-status, def:payment-paid-order-item, def:c08-delivery-status, def:c11-payment-time-limit}
  \notready
  When a payment covering unpaid order items becomes \texttt{SUCCESSFUL},
  every service in those items moves from \texttt{CONFIRMED} to
  \texttt{READY TO PROCEED} and the items lose their payment time limits.
  When it becomes \texttt{FAILED}, the service statuses and payment time
  limits stay as they were.

  Open point: The source states both outcomes for redirected payments only.
  The business must confirm they apply to card and settlement plan payments,
  and decide what a successful payment that covers part of an item's amount
  does to its services.

  Sources: \href{https://standards.atlassian.net/wiki/spaces/EASD/pages/574455927/Pay+Offline+Using+Redirection}{EASD guide, Pay Offline Using Redirection};
  \href{https://standards.atlassian.net/wiki/spaces/EASD/pages/670203914/Payment+Via+Redirection+with+Airline+Notification+of+Payment+Success}{EASD guide, Payment Via Redirection with Airline Notification of Payment Success​};
  \href{https://standards.atlassian.net/wiki/spaces/EASD/pages/574849062/Order}{EASD guide, Order};
  \href{https://standards.atlassian.net/wiki/spaces/EASD/pages/670203914}{EASD guide, Payment Via Redirection with Airline Notification of Payment Success​};
  \href{https://standards.atlassian.net/wiki/spaces/OA/pages/844202130}{Order Accounting, Customer requests a change of a paid order prior to travel resulting in additional collection and a refund}.
\end{theorem}

\begin{theorem}[Time limit on unpaid items (C-129)]
  \label{thm:payment-time-limit-on-unpaid}
  \uses{def:c11-payment-time-limit, def:payment-paid-order-item, def:instant}
  \notready
  Every order item that is unpaid or partly paid carries a payment time
  limit, and a paid item carries none. A time limit of zero on an offer item
  means the item must be paid when the order is created.

  Case to check: An order created without payment returns no time limit for
  one item. Under the second reading that item never expires and holds its
  seat with nothing paid.

  Open point: The ARMI guidance expects the airline to set a limit on every
  unpaid item, while the EASD Time Limits page says a missing limit only
  means no advice was given. The business must decide whether an unpaid item
  may exist without a limit and, if so, whether it can stay unpaid
  indefinitely.

  Sources: \href{https://retailing.iata.org/armi/docs/PAYORD/}{IATA capability PAYORD};
  \href{https://standards.atlassian.net/wiki/spaces/EASD/pages/574423066/Time+Limits}{EASD guide, Time Limits};
  \href{https://standards.atlassian.net/wiki/spaces/EASD/pages/574455879/Create+Order+Without+Payment}{EASD guide, Create Order Without Payment};
  \href{https://standards.atlassian.net/wiki/spaces/EASD/pages/574455927/Pay+Offline+Using+Redirection}{EASD guide, Pay Offline Using Redirection};
  \href{https://retailing.iata.org/armi/docs/ORDWPM/}{IATA capability ORDWPM};
  \href{https://standards.atlassian.net/wiki/spaces/EASD/pages/574849062/Order}{EASD guide, Order}.
\end{theorem}

\begin{definition}[Unpaid item at its time limit (C-130)]
  \label{def:payment-time-limit-expiry}
  \uses{def:c11-payment-time-limit, def:c08-delivery-status, def:c06-item-status, def:payment-payment-status}
  \notready
  When an order item's payment time limit passes without a successful
  payment, the item is removed from the order, and a payment still
  \texttt{PENDING} for it is marked \texttt{FAILED}. The airline may move
  the limit when it accepts a redirected payment as \texttt{PENDING}.

  Case to check: A customer completes a redirected card payment one minute
  after the limit. The airline has already marked the payment
  \texttt{FAILED} and removed the item, so the customer has paid for nothing
  unless the capture is refunded or the item restored.

  Open point: This partly answers the open point of C-11: the Time Limits
  page says removal is expected but names no resulting statuses. The
  business must decide the item status and delivery status after removal
  (for example \texttt{CANCELLED} with \texttt{REMOVED}, or
  \texttt{EXPIRED}) and what happens when a redirected payment is captured
  after the limit.

  Sources: \href{https://standards.atlassian.net/wiki/spaces/EASD/pages/574423066/Time+Limits}{EASD guide, Time Limits};
  \href{https://standards.atlassian.net/wiki/spaces/EASD/pages/574455927/Pay+Offline+Using+Redirection}{EASD guide, Pay Offline Using Redirection};
  \href{https://standards.atlassian.net/wiki/spaces/EASD/pages/1404665857}{EASD guide, 24.1 Release Notes}.
\end{definition}

\begin{theorem}[No payment after the limits (C-131)]
  \label{thm:payment-no-late-payment}
  \uses{def:c11-payment-time-limit, thm:accept-rejects-expired, def:instant}
  \notready
  The airline accepts a payment for an order item only before the item's
  payment time limit and, when the item has a price guarantee, before that
  guarantee expires. A payment attempted after the price guarantee has
  expired fails with a price guarantee error.

  Case to check: The airline extends an item's price guarantee two days past
  its payment time limit, which the source allows. The item can no longer be
  paid while its price is still guaranteed.

  Open point: The payment time limit part is clear, but for the price
  guarantee the source also says repricing is only recommended. The business
  must decide whether an expired guarantee blocks payment or only lets the
  airline change the price, and whether a guarantee may outlast the payment
  time limit, which the source says it should not.

  Sources: \href{https://standards.atlassian.net/wiki/spaces/EASD/pages/574226556/Pay+for+an+Existing+Order}{EASD guide, Pay for an Existing Order};
  \href{https://standards.atlassian.net/wiki/spaces/EASD/pages/574587020/Payment+Errors}{EASD guide, Payment Errors};
  \href{https://standards.atlassian.net/wiki/spaces/EASD/pages/574423066/Time+Limits}{EASD guide, Time Limits}.
\end{theorem}

\begin{theorem}[Failed payment leaves items unpaid (C-132)]
  \label{thm:payment-failed-payment-all-items}
  \uses{def:c13-several-methods, def:payment-paid-order-item}
  One payment transaction may cover several order items with several payment
  methods. If any method's transaction fails, the whole payment fails and
  none of the items it covers becomes paid.

  Case to check: A seller who wants item 1 paid even when the card for item
  2 is declined must send two payments. One combined payment leaves both
  items unpaid.

  Sources: \href{https://standards.atlassian.net/wiki/spaces/EASD/pages/2592899087/Use+of+multiple+payment+methods+to+pay+for+a+single+item}{EASD guide, Use of multiple payment methods to pay for a single item};
  \href{https://standards.atlassian.net/wiki/spaces/EASD/pages/2593325059/Pay+using+multiple+payment+methods+for+a+single+item}{EASD guide, Pay using multiple payment methods for a single item};
  \href{https://standards.atlassian.net/wiki/spaces/EASD/pages/2593030164/Payment+Simulation}{EASD guide, Payment Simulation}.
\end{theorem}

\begin{definition}[Allocating a payment across items (C-133)]
  \label{def:payment-allocation-across-items}
  \uses{thm:payment-failed-payment-all-items, def:payment-paid-order-item, def:c17-refund-methods}
  \notready
  When one payment with several methods covers several order items, the
  airline decides how much of each method goes to each item and the seller
  has no say. A seller who wants a fixed split sends one payment per item.

  Case to check: Items of 800.00 and 100.00 are paid with 400.00 by
  settlement plan and 500.00 by card. If the 100.00 item is cancelled, its
  refund goes to the card under one allocation and to the settlement plan
  under another.

  Open point: The source gives the airline the allocation and states no rule
  for it. The business must fix the rule, because it decides which method is
  refunded when one of the items is later cancelled.

  Sources: \href{https://standards.atlassian.net/wiki/spaces/EASD/pages/2593030164/Payment+Simulation}{EASD guide, Payment Simulation};
  \href{https://standards.atlassian.net/wiki/spaces/EASD/pages/2592899087/Use+of+multiple+payment+methods+to+pay+for+a+single+item}{EASD guide, Use of multiple payment methods to pay for a single item};
  \href{https://retailing.iata.org/armi/docs/PAYMIX/}{IATA capability PAYMIX}.
\end{definition}

\begin{definition}[Combinable payment methods (C-134)]
  \label{def:payment-combinable-methods}
  \uses{def:c13-several-methods}
  \notready
  The payment catalogue places each accepted method in one or more groups,
  and two methods may be combined in one payment when they share a group.
  The airline refuses a payment that combines methods it does not accept
  together.

  Case to check: Method A is in groups 1 and 2, B in groups 2 and 3, and C
  in groups 1 and 3. Every pair shares a group, but no group holds all
  three.

  Open point: The source defines combinability for pairs only. The business
  must decide whether three or more methods need one group that holds them
  all or only a shared group for each pair.

  Sources: \href{https://standards.atlassian.net/wiki/spaces/EASD/pages/2592899087/Use+of+multiple+payment+methods+to+pay+for+a+single+item}{EASD guide, Use of multiple payment methods to pay for a single item};
  \href{https://standards.atlassian.net/wiki/spaces/EASD/pages/2593325059/Pay+using+multiple+payment+methods+for+a+single+item}{EASD guide, Pay using multiple payment methods for a single item}.
\end{definition}

\begin{definition}[Surcharge on each payment method (C-135)]
  \label{def:payment-surcharge-per-method}
  \uses{thm:c14-amount-covers-price, def:c13-several-methods, def:money-add}
  Each method in a payment has a requested amount and a surcharge that
  depends on the method and that amount, and the amount charged to the
  method is the requested amount plus its surcharge. The requested amounts
  of one payment add up to the whole cost of the items it pays.

  Case to check: Reading paid in full as charged amounts equal to the price
  fails on the source's own example: a 100.00 item is paid with 20.00 by
  settlement plan and 82.00 by card, of which 2.00 is the card surcharge.
  This settles most of the open point of C-14.

  Sources: \href{https://standards.atlassian.net/wiki/spaces/EASD/pages/2593030164/Payment+Simulation}{EASD guide, Payment Simulation};
  \href{https://standards.atlassian.net/wiki/spaces/EASD/pages/2592899087/Use+of+multiple+payment+methods+to+pay+for+a+single+item}{EASD guide, Use of multiple payment methods to pay for a single item}.
\end{definition}

\begin{theorem}[Quoted surcharge at payment time (C-136)]
  \label{thm:payment-quoted-surcharge}
  \uses{def:payment-surcharge-per-method, thm:c14-amount-covers-price}
  \notready
  When the airline has quoted a precise surcharge for a stated card brand
  and card type, the payment must use that brand and type and is charged
  that surcharge. A payment with a different brand, scheme, or consumer or
  commercial card type fails with a mismatch error.

  Case to check: The offer is priced with a 3.00 surcharge for a consumer
  Visa debit card, and at payment the airline charges 5.00 for the same
  card. Nothing in the source forbids it.

  Open point: The mismatch errors exist, but no page says a quoted surcharge
  binds the airline, and the catalogue calls its ranges estimates. The
  business must decide whether the charged surcharge may differ from the
  quote or fall outside a published minimum and maximum.

  Sources: \href{https://standards.atlassian.net/wiki/spaces/EASD/pages/574488711/Payment+Catalogue}{EASD guide, Payment Catalogue};
  \href{https://standards.atlassian.net/wiki/spaces/EASD/pages/574685328/Payment+Preference}{EASD guide, Payment Preference};
  \href{https://standards.atlassian.net/wiki/spaces/EASD/pages/574587020/Payment+Errors}{EASD guide, Payment Errors};
  \href{https://standards.atlassian.net/wiki/spaces/ARMI/pages/584318980/General+Guidance}{ARMI, General Guidance}.
\end{theorem}

\begin{definition}[Refundable payment surcharges (C-137)]
  \label{def:payment-surcharge-refundability}
  \uses{def:payment-surcharge-per-method, def:c17-refund-methods, def:accept-refund, thm:refunds-le-payments}
  \notready
  Each payment surcharge carries an indicator saying whether it is
  refundable. A refund of an item's payment returns a surcharge only when
  its indicator marks it refundable.

  Case to check: A 100.00 item is paid with 20.00 by settlement plan and
  82.00 by card, including a non-refundable 2.00 surcharge. A full refund of
  102.00 returns the surcharge, while the cap under the indicator is 100.00.

  Open point: The indicator appears in every simulation example, always
  marked non-refundable, and no page says how a refund uses it. The business
  must decide whether the indicator binds refunds; with refundable
  surcharges a refund can exceed the item's price, which bears on C-17 and
  the open refund question.

  Sources: \href{https://standards.atlassian.net/wiki/spaces/EASD/pages/2593030164/Payment+Simulation}{EASD guide, Payment Simulation};
  \href{https://standards.atlassian.net/wiki/spaces/EASD/pages/2592899087/Use+of+multiple+payment+methods+to+pay+for+a+single+item}{EASD guide, Use of multiple payment methods to pay for a single item}.
\end{definition}

\begin{definition}[Refund as a negative payment (C-138)]
  \label{def:payment-refund-negative-payment}
  \uses{def:refund, def:payment, def:totals, thm:refunds-le-payments, def:c17-refund-methods}
  \notready
  A refund appears in the order as a payment transaction of its own, with a
  new payment ID, a negative amount, a payment method and its own status. In
  the worked examples the method is the one that paid, but the record does
  not point to the payment it returns.

  Case to check: An order holds a 1000.00 settlement plan payment.
  Cancelling it adds a payment of minus 1000.00 with its own ID and no
  reference to the first.

  Open point: The source does not link a refund to the payment it reverses,
  while our model attaches each refund to one payment. The business must
  decide whether to keep that link; without it the per payment limit in the
  open refund question cannot be checked and only an order level limit
  remains.

  Sources: \href{https://retailing.iata.org/armi/docs/PAYREF/}{IATA capability PAYREF};
  \href{https://retailing.iata.org/armi/docs/PAYSUM/}{IATA capability PAYSUM};
  \href{https://standards.atlassian.net/wiki/spaces/EASD/pages/817594417/Customer+requests+a+cancellation+to+paid+Order+or+OrderItems+resulting+in+full+refund}{EASD guide, Customer requests a cancellation to paid Order or OrderItems resulting in full refund};
  \href{https://standards.atlassian.net/wiki/spaces/EASD/pages/842301447/Customer+requests+a+change+of+a+paid+order+prior+to+travel+resulting+in+additional+collection+and+a+refund}{EASD guide, Customer requests a change of a paid order prior to travel resulting in additional collection and a refund}.
\end{definition}

\begin{definition}[Retrying a failed payment (C-139)]
  \label{def:payment-failed-payment-recovery}
  \uses{def:payment-payment-status, thm:payment-payment-outcome-effects, def:c11-payment-time-limit}
  \notready
  After a payment fails, its items stay unpaid with services and time limits
  unchanged, and the seller may try again in a new payment transaction
  before the time limit. The error tells the seller whether to correct the
  card details, add payer authentication, use another method, or stop; after
  a hard stop the airline accepts no payment from that payer.

  Case to check: A payer flagged for suspected fraud on one order pays for a
  second order with the same card. The source does not say whether the stop
  carries over.

  Open point: The source calls recovery optional and does not say whether
  the airline must accept a retry. The business must decide how many retries
  are allowed, whether a hard stop applies to the payer, the order or the
  method, and whether the order is then cancelled.

  Sources: \href{https://standards.atlassian.net/wiki/spaces/EASD/pages/574423154/Payment+Status}{EASD guide, Payment Status};
  \href{https://standards.atlassian.net/wiki/spaces/EASD/pages/574587020/Payment+Errors}{EASD guide, Payment Errors};
  \href{https://standards.atlassian.net/wiki/spaces/EASD/pages/574455927/Pay+Offline+Using+Redirection}{EASD guide, Pay Offline Using Redirection};
  \href{https://standards.atlassian.net/wiki/spaces/EASD/pages/574226627/Payer+Authentication}{EASD guide, Payer Authentication};
  \href{https://retailing.iata.org/armi/docs/PAYRCV/}{IATA capability PAYRCV}.
\end{definition}

\begin{theorem}[Commitment only when nothing due (C-140)]
  \label{thm:payment-commitment-zero-due}
  \uses{def:c12-zero-commitment, def:payment-paid-order-item}
  \notready
  The commitment payment method always carries an amount of zero and may
  cover only items with nothing to pay or collect. Items with nothing due
  may also be covered by a normal payment for a payable item, and every item
  a payment covers, zero priced ones included, is listed on it.

  Open point: The EASD guide added a commitment method in release 26.1,
  while the ARMI guidance and the change worked examples record the same
  commitment as an OT payment of zero. The business must pick one form and
  decide whether a zero OT payment counts when testing whether an item is
  paid.

  Sources: \href{https://standards.atlassian.net/wiki/spaces/EASD/pages/2593488900/Commitment+to+pay+when+no+amount+due}{EASD guide, Commitment to pay when no amount due};
  \href{https://standards.atlassian.net/wiki/spaces/EASD/pages/2533031964/26.1+Release+Notes}{EASD guide, 26.1 Release Notes};
  \href{https://standards.atlassian.net/wiki/spaces/ARMI/pages/584318980/General+Guidance}{ARMI, General Guidance};
  \href{https://standards.atlassian.net/wiki/spaces/EASD/pages/842268768/Customer+requests+a+change+of+a+paid+order+prior+to+travel+resulting+in+a+refund}{EASD guide, Customer requests a change of a paid order prior to travel resulting in a refund}.
\end{theorem}

\begin{definition}[Paying with a voucher (C-141)]
  \label{def:payment-voucher-payment}
  \uses{thm:c20-value-store, def:c17-refund-methods, def:c13-several-methods}
  \notready
  A voucher is a payment method that spends value the airline holds for the
  payer, alone or combined with other methods. Value stored as an order item
  can be spent in another order only after the airline issues it as a
  voucher.

  Case to check: An 800.00 item is paid with a 300.00 voucher and 500.00 by
  card, then cancelled with a 100.00 penalty. Refunding the card first
  returns 500.00 to the card and 200.00 to the voucher; a pro rata split
  returns 437.50 and 262.50.

  Open point: The voucher capability has no guidance beyond its name. The
  business must decide whether a voucher can be partly spent, whether it
  expires, and whether a refund of a voucher paid amount goes back to the
  voucher or to another method.

  Sources: \href{https://retailing.iata.org/armi/docs/PAYVCH/}{IATA capability PAYVCH};
  \href{https://retailing.iata.org/armi/docs/PAYMIX/}{IATA capability PAYMIX};
  \href{https://standards.atlassian.net/wiki/spaces/EASD/pages/574554245/Value+Store}{EASD guide, Value Store};
  \href{https://standards.atlassian.net/wiki/spaces/EASD/pages/574619671/Payment+Concepts}{EASD guide, Payment Concepts}.
\end{definition}
